\section{Related Work}

\paragraph{Sycophancy and Instruction Compliance.}
Instruction-tuned LLMs exhibit a well-documented tendency toward sycophantic behavior: complying with requests and avoiding disagreement even when inaction would be more helpful \cite{perez2023discovering, sharma2024towards, wei2024simple}. OpenAI's April 2025 rollback of a GPT-4o update that made the model excessively agreeable illustrates how deeply this tendency is embedded in RLHF training \cite{openai2025sycophancy}. Our work extends the sycophancy literature from single-turn agreement to multi-turn revision: we show that the same compliance mechanism produces cumulative quality degradation when activated repeatedly.

\paragraph{LLM Self-Refinement.}
\citet{madaan2024self} propose Self-Refine, demonstrating that models can improve outputs through prompted self-critique. \citet{huang2024large} argue that models cannot reliably self-correct reasoning without external feedback, a finding our data supports and extends to the revision setting: undirected revision (without external signal) degrades quality: all six models' quality-optimal stopping point is Turn~1, and the only within-model comparison with sufficient power (Llama, $n = 45$) shows a significant decline ($p < 0.001$). Targeted feedback (specific critique) restores effective revision. \citet{kim2024language} examine conditions under which self-refinement succeeds or fails. These studies typically assume revision is desirable; we examine the complementary question of what happens when models are given repeated opportunities to revise outputs that are already sufficient.

\paragraph{LLM-as-Judge.}
Using LLMs to evaluate generated text has become standard practice \cite{zheng2024judging, li2024generation}. Concerns about self-preferencing bias are well-documented \cite{panickssery2024llm, koo2024benchmarking}. We address this by selecting our evaluator empirically through a six-model judge calibration, choosing the model (Claude Sonnet~4) with highest human-judge correlation ($r = 0.505$), and validating with three human raters (binary threshold agreement 76.6--82.8\% across rater pairs).

\paragraph{LLM Overthinking and Test-Time Scaling.}
\citet{chen2025overthinking} identify a parallel phenomenon in single-turn reasoning: models continue generating past a ``Reasoning Completion Point'' and degrade their own answers through redundant self-correction. \citet{ghosal2025overthinking} demonstrate that test-time scaling in reasoning models produces non-monotonic accuracy curves: accuracy rises with additional thinking tokens up to a critical threshold, then declines steadily as further computation introduces variance and undermines precision. Our multi-turn findings extend this pattern to revision behavior, where the contextual pressure to comply is stronger than in reasoning chains. Together, these results suggest that termination failure is a general property of current LLMs -- inside a single response it manifests as overthinking, and across responses it manifests as our revision cliff. The optimal stopping point exists in both cases, and models systematically overshoot it.

\paragraph{Multi-Turn Dialogue Degradation.}
Multi-turn degradation has been observed in dialogue coherence: \citet{zhang2020dialogpt} show that response quality in open-domain dialogue systems declines over extended conversations, and \citet{thoppilan2022lamda} document coherence decay in long dialogues despite safety filtering. However, these studies focus on conversational coherence rather than revision quality. Our work quantifies a distinct phenomenon: degradation caused not by topic drift in dialogue but by the revision act itself, where models regenerate outputs at progressively lower quality even when the task specification remains fixed.

\paragraph{Self-Correction Without External Feedback.}
\citet{kamoi2024selfcorrection} provide a critical survey of LLM self-correction in TACL, concluding that intrinsic self-correction (without external feedback) generally degrades performance across tasks. Our targeted-feedback result is a quantified confirmation in the revision setting: undirected revision degrades quality by 0.74 levels (after stripping meta-commentary, $n = 50$) while targeted feedback improves by 1.16 levels ($p < 0.0001$, $n = 177$), consistent with their conclusion that reliable external signal is what makes self-correction work. \citet{shinn2023reflexion} propose Reflexion, where agents improve through verbal self-reflection stored in episodic memory; notably, they identify a ``degeneration-of-thought'' failure mode where repeated self-reflection without new information converges on worse outputs -- a conceptual precursor to our revision cliff in the multi-turn conversation setting.

\paragraph{Reward Hacking and Verbosity Bias.}
Reward hacking occurs when models exploit proxy objectives rather than the intended goal \cite{skalse2022defining}. \citet{singhal2023long} document verbosity bias in RLHF-trained models, where longer outputs receive higher reward scores regardless of quality, incentivizing padding. Our evaluator design mitigates verbosity bias through scale anchoring: the five-level rubric (six levels with a 6$\to$2 recode for non-monotonic ``Overdone'') explicitly penalizes content loss and rewards task adherence rather than length. The revision cliff may represent a novel form of reward hacking where the proxy signal is visible effort (revising) rather than output length.

\paragraph{Token Efficiency and AI Cost.}
Enterprise LLM API spend roughly doubled from \$3.5B to \$8.4B in under a year \cite{menlo2025llmmarket}. Per-token prices have fallen sharply since 2023 (roughly 80\% by some industry estimates), yet total enterprise AI spend has grown several-fold over the same period because volume outran price declines, a pattern that makes per-token waste economically significant even as unit costs shrink. \citet{deloitte2026stateofai} and \citet{mckinsey2025stateofai} document growing concern about AI deployment costs relative to realized value. Our revision tax quantifies a specific, measurable mechanism contributing to this waste: undirected revision that the model cannot decline, producing tokens that make outputs worse while increasing costs.

\paragraph{Multi-Turn Task Performance.}
Closest to our work is \citet{laban2025lost}, who demonstrate that LLMs lose 39\% accuracy in multi-turn conversations compared to single-turn equivalents, earning the ICLR 2026 Best Paper Award. Their degradation mechanism is distinct from ours: it arises from underspecification, where task requirements are revealed gradually across turns and models make premature assumptions in early turns that they cannot recover from. In contrast, our tasks are fully specified from the outset and the output is already sufficient -- degradation is caused by the revision act itself rather than by information asymmetry. The two findings are complementary: theirs shows models fail when requirements accumulate across turns; ours shows models fail when asked to revise outputs that already satisfy all requirements.
